% Descriptive Statistics Section - LaTeX Code

\section{Descriptive Statistics}

\subsection{Sample Characteristics}

The final analytical sample consisted of 146 participants who completed a total of 6,800 recognition memory trials, averaging approximately 46.6 trials per participant. Participants were distributed across two experimental conditions: 79 participants (54.1\%) were assigned to the Natural Boundary (NB) condition, while 67 participants (45.9\%) were assigned to the Abrupt Boundary (AB) condition. 

The sample was predominantly young adults with a mean age of 22.0 years (\textit{SD} = 1.9 years, range: 19--28 years, \textit{Mdn} = 22 years). Age distributions were comparable across conditions, with NB participants averaging 22.2 years (\textit{SD} = 2.1 years) and AB participants averaging 21.9 years (\textit{SD} = 1.6 years). The sample was predominantly male (112 males, 76.7\%; 34 females, 23.3\%), with similar gender distributions across both conditions (NB: 63 males, 16 females; AB: 49 males, 18 females). Nearly all participants were right-handed (139, 95.2\%), with only 7 left-handed individuals (4.8\%), consistent with the general population distribution. All participants reported either normal vision (100 participants, 68.5\%) or corrected-to-normal vision (46 participants, 31.5\%), ensuring that visual perception did not confound the results.

\subsection{Recognition Accuracy}

Overall recognition accuracy across all participants was high, with a mean accuracy of 0.857 (\textit{SD} = 0.077, or 85.7\% $\pm$ 7.7\%), indicating that participants correctly identified the target frame in approximately 86\% of trials. The median accuracy was 0.875 (87.5\%), and performance ranged from 0.575 (57.5\%) to 0.988 (98.8\%). The interquartile range (IQR) spanned from 0.825 to 0.900, indicating that the middle 50\% of participants scored between 82.5\% and 90\% accuracy, reflecting relatively consistent performance across the sample.

When examined by experimental condition, a notable difference emerged. Participants in the NB condition demonstrated higher recognition accuracy (\textit{M} = 0.872, \textit{SD} = 0.069, \textit{Mdn} = 0.875) compared to those in the AB condition (\textit{M} = 0.838, \textit{SD} = 0.084, \textit{Mdn} = 0.850), representing a 3.4 percentage point difference favoring the NB group. This corresponds to approximately a 4.0\% relative improvement in recognition accuracy when event boundaries were preserved. The AB condition exhibited greater performance variability (\textit{SD} = 0.084) than the NB condition (\textit{SD} = 0.069), suggesting that abrupt boundary disruptions produced less consistent encoding quality across participants. Furthermore, the AB condition included participants with substantially lower performance (minimum = 0.575 or 57.5\%) compared to the NB condition (minimum = 0.675 or 67.5\%), while peak performance was similar across conditions (AB maximum = 0.975; NB maximum = 0.988).

\subsection{Response Time}

Mean response time across all participants was 5.718 seconds (\textit{SD} = 1.571 seconds), with a median response time of 5.470 seconds, indicating that participants typically required between 4.7 and 6.2 seconds (IQR: 4.740--6.158 seconds) to make their recognition decisions. Response times ranged from 3.251 seconds to 12.180 seconds, reflecting substantial individual differences in decision speed.

Contrary to a simple speed-accuracy tradeoff, participants in the NB condition exhibited longer mean response times (\textit{M} = 5.834 seconds, \textit{SD} = 1.636 seconds, \textit{Mdn} = 5.534 seconds) compared to those in the AB condition (\textit{M} = 5.582 seconds, \textit{SD} = 1.491 seconds, \textit{Mdn} = 5.399 seconds). This 0.25-second difference suggests that NB participants took slightly more time to respond yet achieved higher accuracy, indicating that they were not simply sacrificing accuracy for speed. This pattern is consistent with the hypothesis that preserved event boundaries facilitate more deliberate retrieval of richer memory traces, whereas disrupted boundaries may lead to faster but less accurate decisions, potentially involving more guessing or reliance on incomplete memory representations.

\subsection{Confidence Ratings}

Participants' confidence in their recognition judgments was assessed on a 5-point Likert scale (1 = \textit{very unconfident}, 5 = \textit{very confident}). Overall, participants reported relatively high confidence in their responses, with a mean confidence rating of 4.149 (\textit{SD} = 0.422, \textit{Mdn} = 4.225), corresponding to approximately 83\% of the maximum possible confidence level. Confidence ratings ranged from 2.900 to 5.000, with the middle 50\% of participants reporting confidence between 3.825 and 4.425 (IQR).

Confidence ratings differed between experimental conditions in a pattern consistent with accuracy differences. NB participants reported higher mean confidence (\textit{M} = 4.209, \textit{SD} = 0.407, \textit{Mdn} = 4.275, range: 3.000--5.000) compared to AB participants (\textit{M} = 4.077, \textit{SD} = 0.431, \textit{Mdn} = 4.100, range: 2.900--4.900). This 0.132-point difference suggests that participants in the NB condition not only performed more accurately but also possessed greater subjective certainty in their memory judgments. The alignment between objective accuracy and subjective confidence indicates adequate metacognitive awareness, with participants able to monitor the quality of their memory retrieval. Additionally, the AB condition showed slightly greater variability in confidence ratings (\textit{SD} = 0.431) compared to the NB condition (\textit{SD} = 0.407), mirroring the pattern observed in accuracy variability.

\subsection{Trial-Level Performance}

At the trial level, 5,822 of the 6,800 trials (85.6\%) resulted in correct target identification, while 978 trials (14.4\%) resulted in incorrect responses. This trial-level accuracy closely matched the participant-level aggregate accuracy, providing convergent evidence for the observed performance patterns. When analyzed by condition, AB trials yielded 84.0\% accuracy (2,720 correct out of 3,240 total trials), whereas NB trials yielded 87.1\% accuracy (3,102 correct out of 3,560 total trials). This 3.1 percentage point difference at the trial level corroborates the participant-level findings, indicating that the NB advantage is robust across both levels of analysis.

Response choice distribution across the entire sample was balanced, with participants selecting the right-side frame (``r'' response) on 3,407 trials (50.1\%) and the left-side frame (``l'' response) on 3,393 trials (49.9\%). This near-equal split indicates that there was no systematic spatial response bias, confirming that the left-right positioning of target and lure frames was adequately counterbalanced across trials.

% Optional: Add a summary paragraph
\subsection{Summary}

In summary, the descriptive statistics reveal several key patterns in the data. First, participants in the NB condition demonstrated consistently superior recognition memory performance compared to those in the AB condition, with higher accuracy (87.2\% vs.\ 83.8\%), higher confidence (4.21 vs.\ 4.08), and lower performance variability. Second, the AB condition was associated with faster but less accurate responses, suggesting possible differences in retrieval strategy rather than a simple speed-accuracy tradeoff. Third, the alignment between accuracy and confidence across conditions indicates preserved metacognitive monitoring, suggesting that the experimental manipulation affected memory strength without impairing participants' ability to assess their own memory quality. These patterns provide preliminary support for the hypothesis that natural event boundaries facilitate more robust and consistent memory encoding compared to artificially disrupted boundaries.
